\section{Introduction}

\subsection{Motivation}

Standard transformer attention (Vaswani et al., 2017) computes:

\[
\text{Attn}(Q, K, V) = \text{softmax}\!\left(\frac{Q K^\top}{\sqrt{d_h}}\right) V
\tag{1.1}
\]

This scales as $\mathcal{O}(L^2)$ in both time and memory, making 
long-context inference expensive.

\textbf{Linear attention} replaces softmax with a kernel feature map:

\[
\text{Attn}(Q, K, V) \approx \frac{\phi(Q) \left(\phi(K)^\top V\right)}{\phi(Q) \left(\phi(K)^\top \mathbf{1}\right)}
\tag{1.2}
\]

Since $\phi(K)^\top V \in \mathbb{R}^{F \times d_h}$ is computed once and 
reused, the total complexity is $\mathcal{O}(L F d_h)$ -- \textbf{linear in L}.

The quality of linear attention depends entirely on how well $\phi$ 
approximates the softmax kernel:

\[
\kappa(q, k) = \exp(q^\top k)
\tag{1.3}
\]

\subsection{The Kernel Approximation Problem}

Three families of feature maps have been studied:

\begin{enumerate}
\item \textbf{Performer (Choromanski et al., 2021):} Positive Orthogonal Random
  Features (PORF), which approximate $\exp(q^\top k)$ with positive
  features.
\item \textbf{Linear Transformer (Katharopoulos et al., 2020):} $\phi(x) = \text{elu}(x) + 1$,
  which corresponds to a specific kernel that is \textbf{not} the softmax
  kernel.
\item \textbf{cosFormer (Qin et al., 2022):} uses $\text{ReLU}$ with a cosine
  reweighting.

\end{enumerate}
We take a fourth approach: \textbf{Random Fourier Features (RFF)} for the 
\textbf{Gaussian RBF kernel} $\exp(-\|q - k\|^2 / 2)$:

\[
\kappa(q, k) \approx \phi(q)^\top \phi(k), \quad
\phi(x) = \frac{1}{\sqrt{F}} \left[ \cos(x W), \sin(x W) \right]
\tag{1.4}
\]

with $W_{ij} \sim \mathcal{N}(0, 1)$. This is the \textbf{real part} of the 
classical RFF estimator (Rahimi \& Recht, 2007).

\subsection{Our Approach: LEKAv1}

LEKAv1 has four components:

\subsubsection{Head-wise LayerNorm}

Before projection, we normalize each head:

\[
\tilde{x} = \text{LayerNorm}(x) \in \mathbb{R}^{d_h}
\tag{1.5}
\]

This centers and scales the input, making the random projection more 
stable.

\subsubsection{QR-Orthogonalized Projection}

The projection matrix $W \in \mathbb{R}^{d_h \times F}$ is initialized from 
the QR decomposition of a random Gaussian:

\[
W = Q, \quad \text{QR}(G) = QR, \quad G \sim \mathcal{N}(0, I)
\tag{1.6}
\]

This ensures orthonormal columns, reducing variance in the kernel 
approximation.

\subsubsection{cos/sin Feature Map}

\[
\phi(x) = \frac{1}{\sqrt{F}} [\cos(\tilde{x} W) \| \sin(\tilde{x} W)]
\tag{1.7}
\]

which gives $\phi(x) \in \mathbb{R}^{2F}$.

\subsubsection{Preservation Gate}

A single learnable scalar $\alpha = \sigma(w)$ blends attention output 
with the input:

\[
\text{out} = \alpha \odot A + (1 - \alpha) \odot x
\tag{1.8}
\]

This provides a \textbf{residual safety net}: if the kernel approximation is 
poor, the model can lean toward the identity.

\subsection{Contributions}

\begin{itemize}
\item \textbf{A production-hardened RFF attention.} We implement cos/sin feature
  maps with QR-orthogonalized projections, causal cumsum, and
  bidirectional global KV.
\item \textbf{Head-wise LayerNorm before projection.} We normalize the head input
  before random projection, which stabilizes the kernel approximation.
\item \textbf{A preservation gate.} A scalar gate controls mixing with the residual.
\item \textbf{Multi-architecture support.} DistilBERT, RoBERTa, GPT-2, Qwen3.5
  with GQA and RoPE.
\item \textbf{An honest evaluation.} We report both the accuracy and the parameter
  cost.

\end{itemize}
\subsection{Paper Organization}

Section 2 reviews related work. Section 3 provides background on kernels, 
RFF, and linear attention. Section 4 presents LEKAv1. Section 5 describes 
integration. Section 6 details the setup. Section 7 reports results. 
Section 8 analyzes the kernel approximation. Section 9 discusses guidance. 
Section 10 lists limitations. Section 11 concludes.



\subsection{Notation}

\begin{table}[H]
\centering
\small
\begin{tabular}{ll}
\toprule
\textbf{Symbol} & \textbf{Meaning} \\
\midrule
$L$ & Sequence length \\
$d$ & Hidden size \\
$H$ & Heads \\
$d_h$ & Head dim \\
$F$ & Feature dimension \\
$\phi$ & Kernel feature map \\
$W$ & Random projection \\
$Q, K, V$ & Attention projections \\
$\alpha$ & Preservation gate \\
\bottomrule
\end{tabular}
\end{table}

